% ==================================================================
% id:       lim-011
% topic:    גבולות ורציפות של פונקציות
% subtopic: גבולות אינסופיים ואסימפטוטות
% solution: full
% ==================================================================
\textbf{
בשאלה זו נקצר ונכתוב
$\ds{\lim_{x\to x_0}}f(x)=\pm\infty$
אם
$\ds{\lim_{x\to x_0}}f(x)=\infty$
או
$\ds{\lim_{x\to x_0}}f(x)=-\infty$.
\\
הוכיחו או הפריכו
כל אחד מהסעיפים הבאים:
}
\begin{enumerate}
    \item אם
    $\ds{\lim_{x\to \infty}}f(x)=\infty$ ו-
    $g(x)$
    חסומה,
    אז
    $\ds{\lim_{x\to \infty}}(f(x)+g(x))=\infty$.

\ifsol
\hint{בדקו אם החסם התחתון של $g$ מאפשר להשוות את $f+g$ לפונקציה ששואפת לאינסוף; אם נראה לכם שהטענה לא נכונה, חפשו דוגמה נגדית פשוטה.}
\begin{solbox}
    הטענה נכונה.
    קיים
    $m$
    כך ש-
    $g(x)>m$
    לכל
    $x\in\mathbb{R}$.
    ולכן
    $$,f(x)+g(x)>f(x)+m$$
    לכל
    $x\in\rr$,
    והרי מאריתמטיקה במובן הרחב ידוע כי
    $$\lim_{x\to \infty} (f(x)+m)=\infty$$
    לכן מכלל הפיצה נסיק שמתקיים
$\lim_{x\to\infty}(f(x)+g(x))=\infty$.
\end{solbox}
\fi

    \item אם
    $\ds{\lim_{x\to 3}}f(x)=\infty$ ו-
    $g(x)$
    חסומה,
    אז
    $\ds{\lim_{x\to 3}}f(x)\cdot g(x)=\pm\infty$.

\ifsol
\hint{חשבו מה קורה למכפלה כש-$g$ מקבלת ערכים מסוגים שונים ליד $3$. אם נראה לכם שהטענה לא נכונה, חפשו דוגמה נגדית פשוטה; אחרת נסו להוכיח אותה.}
\begin{solbox}
הטענה אינה נכונה.
נביא דוגמא נגדית:
$$.f(x)=\frac{1}{|x-3|},\ g(x)=\case{cc}{|x-3| & x\in[2,4]\\ 1 & x\notin[2,4]} $$
מתקיים
$\lim_{x\to3}f(x)=\infty$,
$|g(x)|\leq1$,
כלומר
$g(x)$
חסומה,
אבל
$$.\lim_{x\to3}f(x)\cdot g(x)=1$$
גם הפונקציה הקבועה
$g(x)=0$
תעבוד במקרה זה כדוגמה נגדית.
\end{solbox}
\fi

    \item אם
    $\ds{\lim_{x\to -2}}f(x)=\infty$ ו-
    $11\leq g(x)\leq13$
    אז
    $\ds{\lim_{x\to -2}}f(x)\cdot g(x)=\pm\infty$.

\ifsol
\hint{חשבו מה החסם התחתון החיובי של $g$ נותן למכפלה $f\cdot g$ ליד $-2$; אם נראה לכם שהטענה לא נכונה, חפשו דוגמה נגדית פשוטה.}
\begin{solbox}
הטענה נכונה.
נקבע
חיובי
$M$,
אזי מכיוון ש-
$\displaystyle\lim_{x\to-2}f(x)=\infty$,
קיים
$\delta>0$
כך ש-
$f(x)>\frac{M}{11}$
לכל
$x\in(-2-\delta,-2+\delta)$.
מתקיים
$$,f(x)\cdot g(x)>\frac{M}{11}\cdot11=M,\ \forall x\in(-2-\delta,2+\delta)$$
משמעו
$$.\ds{\lim_{x\to -2}}f(x)\cdot g(x)=\infty$$
\end{solbox}
\fi

    \item אם
    $\ds{\lim_{x\to \infty}}f(x)=\infty$ ו-
    $g(x)$
    חסומה וחיובית,
    אז
    $\ds{\lim_{x\to \infty}}f(x)\cdot g(x)=\infty$.

\ifsol
\hint{השוו לסעיף הקודם: האם חיוביות של $g$ נותנת את מה שהחסם התחתון החיובי נתן שם? אם נראה לכם שהטענה לא נכונה, חפשו דוגמה נגדית פשוטה.}
\begin{solbox}
הטענה אינה נכונה.
נביא דוגמא נגדית:
$$.f(x)=x^2+1,\ g(x)=\frac{1}{x^2+1}$$
מתקיים
$\displaystyle\lim_{x\to\infty}f(x)=\infty$,
$g(x)\leq 1$,
כלומר,
$g(x)$
חסומה, אבל
$$.\lim_{x\to\infty}f(x)\cdot g(x)=1$$
\end{solbox}
\fi

    \item אם
    $\ds{\lim_{x\to\infty}}f(x)\cdot g(x)=\infty$,
    אז
    $\ds{\lim_{x\to\infty}}f(x)=\pm\infty$
    ו-
    $\ds{\lim_{x\to\infty}}g(x)=\pm\infty$.

\ifsol
\hint{נסו מכפלות של חזקות של $x$ ששואפות לאינסוף ובדקו את הגבול של כל גורם בנפרד; אם הטענה נכונה בכל הדוגמאות, נסו להוכיח אותה.}
\begin{solbox}
    הטענה אינה נכונה.
    נציג דוגמא נגדית:
    $$.f(x)=\frac{1}{x},\ g(x)=x^2$$
    מתקיים
    $$,\lim_{x\to\infty}f(x)\cdot g(x)=\lim_{x\to\infty}x=\infty$$
    אבל
    $\displaystyle\lim_{x\to\infty}f(x)=0\neq\pm\infty$.
\end{solbox}
\fi

    \item אם
    $\ds{\lim_{x\to\infty}}f(x)\cdot g(x)=\infty$,
    אז
    $\ds{\lim_{x\to\infty}}f(x)=\pm\infty$
    או
    $\ds{\lim_{x\to\infty}}g(x)=\pm\infty$.

\ifsol
\hint{השוו לסעיף הקודם. אם נראה לכם שהטענה לא נכונה, חפשו דוגמה נגדית פשוטה; אחרת נסו להוכיח אותה.}
\begin{solbox}
        הטענה אינה נכונה.
    נציג דוגמא נגדית:
    $$.f(x)=g(x)=(-1)^{\lfloor x\rfloor}\cdot x$$
    מתקיים
    $$,\lim_{x\to\infty}f(x)\cdot g(x)=\lim_{x\to\infty}x^2=\infty$$
    אבל
    $\displaystyle\lim_{x\to\infty}f(x)\neq\pm\infty$,
    למעשה, הגבול לא קיים כלל.
    אכן,
    $$,f(x)=x,\quad \forall\ 2k\leq x<2k+1,\ k\in\mathbb{Z}$$
    מצד שני
    $$.f(x)=-x,\quad \forall\ 2k+1\leq x<2(k+1),\ k\in\mathbb{Z}$$
\end{solbox}
\fi

    \item אם
    $\ds{\lim_{x\to1}}\frac{f(x)}{g(x)}=\infty$,
    אז
    $\ds{\lim_{x\to1}}(f(x)-g(x))=\infty$.

\ifsol
\hint{חשבו באילו דרכים מנה יכולה לשאוף לאינסוף, ומה כל אחת מהן אומרת על ההפרש. אם נראה לכם שהטענה לא נכונה, חפשו דוגמה נגדית פשוטה; אחרת נסו להוכיח אותה.}
\begin{solbox}
    הטענה אינה נכונה. להלן דוגמא נגדית:
    $$,f(x)=1,\ g(x)=|x-1|$$
    מתקיים
    $$,\lim_{x\to1}\frac{f(x)}{g(x)}=\lim_{x\to1}\frac{1}{|x-1|}=\infty$$
    אבל
    $$.\lim_{x\to1}(f(x)-g(x))=\lim_{x\to1}(1-|x-1|)=1\neq\infty$$
\end{solbox}
\fi

    \item אם
    $\ds{\lim_{x\to-5}}\frac{f(x)}{g(x)}=\infty$,
    ו-
    $\ds{\lim_{x\to-5}}g(x)=\infty$,
    אז
    $\ds{\lim_{x\to-5}}(f(x)-g(x))=\infty$.

\ifsol
\hint{נסו להוציא את $g$ כגורם משותף מ-$f-g$ ובדקו לאן שואף כל גורם; אם זה לא מוביל לתשובה, חפשו דוגמה נגדית.}
\begin{solbox}
    הטענה נכונה.
    נשים לב כי מתקיים
    $$f(x)-g(x)=g(x)\cdot\left(
    \frac{f(x)}{g(x)}-1
    \right)$$
    ולכן מאריתמטיקה במובן הרחב של חיסור וכפל נקבל כי
    $$.\lim_{x\to-5}
    (f(x)-g(x))=\infty$$
\end{solbox}
\fi
\end{enumerate}
